% from ega0-21-fr.tex
\oldpage[0\textsubscript{IV-1}]{154}
\subsection{Systèmes de $p$-générateurs et $p$-bases}
\label{0.21.1-fr}

\begin{env}[21.1.1]
\label{0.21.1.1-fr}
Étant donné un nombre $p$ qui est, soit $0$, soit un nombre premier, nous
dirons qu'un anneau $A$ est de caractéristique $p$ s'il existe un homomorphisme
d'anneaux $P\to A$, où $P$ est le corps premier de caractéristique $p$~; on
notera que cet homomorphisme est alors unique, le composé
$\mathbf Z\to P\to A$ étant l'unique homomorphisme de $\mathbf Z$ dans $A$.
Si $A\neq0$, il revient au même de dire que $A$ contient un corps de
caractéristique $p$, l'image de $P$ étant nécessairement un corps isomorphe à
$P$ (et d'ailleurs le seul sous-corps de $A$ isomorphe à $P$).
\end{env}

\begin{env}[21.1.2]
\label{0.21.1.2-fr}
Si $p>0$, dire que $A$ est de caractéristique $p$ équivaut à dire que, dans
$A$, on a $p\cdot1=0$, ou encore $pA=0$. Si $p=0$, dire que $A$ est de
caractéristique $p$ équivaut à dire que pour tout entier $n\neq0$, $n\cdot1$
est inversible dans $A$. Si $A\neq0$, il ne peut exister qu'un seul $p$
(premier ou $0$) tel que $A$ soit de caractéristique $p$~; cela résulte de
l'identité de Bézout $ap+bq=1$ pour deux nombres premiers distincts $p,q$.
Par contre, l'anneau réduit à $0$ est de caractéristique $p$ pour tout $p$.
\end{env}

\begin{env}[21.1.3]
\label{0.21.1.3-fr}
Si $A$ est de caractéristique $p$, il en est de même de toute algèbre sur
$A$. En particulier, pour tout idéal premier $\mathfrak p$ de $A$, le corps
résiduel de $A$ en $\mathfrak p$ est de caractéristique $p$. Réciproquement,
si $p=0$ et si, pour tout idéal maximal $\mathfrak m$ de $A$, le corps
résiduel de $A$ en $\mathfrak m$ est de caractéristique $0$, il en est de même
de $A$, car pour tout entier $n\neq0$, $n\cdot1$ est alors inversible dans
tous les $A_{\mathfrak m}$, donc aussi dans $A$. Par contre, si $p>0$, un
anneau local peut avoir son corps résiduel de caractéristique $p$ sans être
lui-même de caractéristique $p$, comme le montre l'exemple de l'anneau
premier $\mathbf Z_{p\mathbf Z}$ (19.8.3) ou de l'anneau local artinien
$\mathbf Z/p^n\mathbf Z$ pour $n\geq2$, qui ne contiennent pas de corps.

Notons enfin que, pour un anneau (même réduit), les corps résiduels en ses
idéaux premiers peuvent avoir des caractéristiques différentes, comme le
montre l'exemple de $\mathbf Z$.
\end{env}

\begin{env}[21.1.4]
\label{0.21.1.4-fr}
Dans toute la suite de ce paragraphe, nous supposons fixé un nombre premier
$p$ et tous les anneaux seront supposés de caractéristique $p$, sauf mention
expresse du contraire. Pour un tel anneau $A$, l'application $x\mapsto x^p$
est un endomorphisme de $A$, que l'on note $F_A$~; si $A$ est réduit, $F_A$
est injectif. On pose $A^p=F_A(A)$ (sous-anneau de $A$ formé des $x^p$,
$x\in A$)~; on peut naturellement considérer $A$ comme une $A^p$-algèbre.

On peut aussi considérer $A$ comme une $A$-algèbre au moyen de
l'homomorphisme $F_A:x\mapsto x^p$ de $A$ dans $A$~; autrement dit, il
s'agit de l'algèbre $A^{(p)}$ pour laquelle le produit $\lambda\cdot x$ de
$x\in A$ par un scalaire $\lambda\in A$ est le produit $\lambda^p x$ dans
l'anneau $A$. Il est clair que, pour tout homomorphisme d'anneaux $u:A\to B$,
le couple $(u,u)$ est un di-homomorphisme d'algèbres $A^{(p)}\to B^{(p)}$.
Pour tout $A$-module $E$, on pose $E^{(p)}=E\otimes_A A^{(p)}$, où
$A^{(p)}$ est considéré comme un $(A,A)$-bimodule, la structure de
$A$-module à gauche étant celle qu'on vient de définir et la structure de
$A$-module à droite définie par la multiplication dans $A$~; ainsi, pour
$x\in E$, $\alpha,\beta\in A$, on a
\[
  \alpha(x\otimes\beta)=x\otimes(\alpha\beta),\qquad
  (\alpha x)\otimes\beta=x\otimes\alpha^p\beta.
\]
\end{env}

